%! Unit Opener: Why Algebra Fluency Matters — Unit 1, Lesson 1.0 — Guided Notes (Answer Key)
\documentclass[10pt]{article}
\usepackage{atda-article}
\usepackage{atda-key}   % pulls in -boxes; do NOT also load -boxes
\newcommand{\TallMath}[1]{$\displaystyle #1\rule[-1.4em]{0pt}{3.2em}$}
\setlength{\workrowsep}{6pt}

\begin{document}

\pageheader{Unit 1, Lesson 1.0}{Guided Notes --- Answer Key}
% NAMESTRIP: no \namedateperiod here — mirrors the blank. Teacher-only prose goes
% in the lesson plan as \begin{teachernote}[Guided Notes], never in a key.

\vspace{0.15in}

\begin{objectivebox}
\small
By the end of this lesson, I will be able to\dots
\begin{enumerate}[leftmargin=1.5em, itemsep=2pt, topsep=3pt]
  \item explain what the \textbf{expanded} form of a model shows me and what the
        \textbf{factored} form shows me, and choose the one that answers a given question;
  \item show that two expressions are \textbf{equivalent} by expanding one and comparing it,
        term by term, with the other;
  \item use the algebra Unit 1 is built on: the exponent rules, the distributive property,
        greatest common factors, and solving a linear equation;
  \item name what each lesson of Unit 1 teaches and where that skill gets used later in this
        course.
\end{enumerate}
\end{objectivebox}

\vspace{0.10in}

\boxguard
\begin{vocabbox}
\small
Fill in each term as we build it together. These five words carry the whole unit.
\vocabans{Equivalent expressions}{Two expressions that give the same value for every allowed
input. You can rewrite one into the other without changing what it means.}
\vocabans{Polynomial}{A sum of terms whose variable exponents are whole numbers, such as
$-16t^2+48t+64$. The degree is the largest exponent.}
\vocabans{Expanded form}{A polynomial written as a sum of terms. It shows the constant term ---
the starting value of the model --- at a glance.}
\vocabans{Factored form}{The same polynomial written as a product of factors. It shows the zeros
at a glance, one from each factor.}
\vocabans{Zero of an expression}{An input that makes the expression equal $0$. In context a zero
often marks when something lands, runs out, or breaks even.}
\end{vocabbox}

\vspace{0.10in}

\boxguard
\begin{hookbox}
\small
\textbf{The Fourth of July shell.} A fireworks shell is launched from a platform on a $64$-foot
cliff above the harbor. Its height above the water, in feet, $t$ seconds after launch is
\[ h(t) = -16t^2+48t+64 \qquad\text{and also}\qquad h(t) = -16(t-4)(t+1). \]
The safety officer needs to know \textbf{when the shell hits the water}. The announcer needs to
know \textbf{how high the platform is}. Which form would you hand to each person, and why?

\ansline{Safety officer: the factored form --- its zero $t=4$ is the landing time.}\\
\ansline{Announcer: the expanded form --- its constant $64$ is the launch height.}\\
\end{hookbox}

\vspace{0.10in}

\boxguard[14]
\begin{notesbox}{1. Are They Really the Same?}
\small
Before we trust either form, check them against each other. Expand the factored form and compare
it, term by term, with the expanded form. This is the check that will verify every factoring
answer you write this year.
\begin{work}
  -16(t-4)(t+1) &= -16\big(t^2+t-4t-4\big) \\
                &= -16\big(t^2-3t-4\big) \\
                &= -16t^2+48t+64
\end{work}

\textbf{Careful:} the $-16$ must reach \ans{every} term inside the parentheses, not just the
first one. Two expressions are \textbf{equivalent} only when \emph{every} coefficient matches.
\end{notesbox}

\vspace{0.10in}

\boxguard[20]
\begin{notesbox}{2. Two Forms, Two Questions}
\small
The two forms give the same output for every $t$. They differ only in what they let you see
without doing any work.

\renewcommand{\arraystretch}{1.5}
\begin{tabularx}{\linewidth}{@{}p{5.3cm} c X@{}}
\rowcolor{mist}
\textbf{The question} & \textbf{Form that answers it} & \textbf{What you read off} \\
\hline
How high is the launch platform? & \ans{expanded} & \ans{the constant, $64$ feet} \\
When does the shell hit the water? & \ans{factored} & \ans{the zero $t=4$ seconds} \\
\end{tabularx}

\vspace{6pt}
The factored form gives two zeros, $t=4$ and $t=-1$. Both are real numbers, but only one of them
means anything for this shell. Which one, and why?

\ansline{$t=4$ --- the shell hits the water 4 seconds after launch.}\\
\ansline{$t=-1$ is one second before the launch, outside the situation.}\\
\end{notesbox}

\vspace{0.10in}

\boxguard[26]
\begin{notesbox}{3. The Unit 1 Roadmap}
\small
Mark the last column for yourself using your warm-up results: \textbf{C} = confident,
\textbf{R} = rusty, \textbf{N} = new to me. You will come back to this page before the unit test.

\renewcommand{\arraystretch}{1.45}
\begin{tabularx}{\linewidth}{@{}c p{2.6cm} X c@{}}
\rowcolor{mist}
\textbf{Lesson} & \textbf{Standard} & \textbf{What I will be able to do} & \textbf{C/R/N} \\
\hline
1.1 & \texttt{A2.EO.3a} & Simplify expressions; add, subtract, and multiply polynomials using
      the exponent rules & \blank{0.9cm} \\
1.2 & \texttt{A2.EO.2a--c} & Simplify radicals and move between radicals and rational exponents
      & \blank{0.9cm} \\
1.3 & \texttt{A2.EO.3b} & Factor out a greatest common factor and factor a four-term polynomial
      by grouping & \blank{0.9cm} \\
1.4 & \texttt{A2.EO.3b} & Factor trinomials, with and without a leading coefficient
      & \blank{0.9cm} \\
1.5 & \texttt{A2.EO.3b, d} & Factor and verify special forms: difference of squares, sum and
      difference of cubes, perfect square trinomials & \blank{0.9cm} \\
1.6 & \texttt{A2.EI.2b}, \texttt{A2.EI.6a--b} & Solve polynomial equations by factoring, and
      interpret the solutions in context & \blank{0.9cm} \\
1.7 & \texttt{A2.EO.1a--b} & Simplify rational expressions by factoring and cancelling
      & \blank{0.9cm} \\
\end{tabularx}

\vspace{4pt}
\textbf{My goal for this unit:} \ansline{student response --- accept any specific, checkable goal}
\end{notesbox}

\vspace{0.10in}

\boxguard
\begin{notesbox}{4. Where This Algebra Gets Spent}
\small
None of this unit is algebra for its own sake. Every skill above is cashed in later this year:

\begin{itemize}[leftmargin=1.3em, itemsep=3pt, topsep=3pt]
  \item The \textbf{exponent rules} of Lesson 1.1 become the \ans{logarithm} properties in
        Unit 3, where $\log(ab)=\log a+\log b$ is the same idea written backwards.
  \item \textbf{Factoring} (1.3--1.5) is what makes \ans{polynomial} equations solvable in
        Lesson 1.6 --- and every trigonometric equation you meet in Unit 8 ends the same way.
  \item \textbf{Simplifying quotients} (1.7) is exactly the move that reduces the counting and
        probability formulas in Unit 4 to numbers you can compute by hand.
\end{itemize}
\end{notesbox}

\vspace{0.10in}

\boxguard[22]
\begin{practicebox}
\small
A gardener has $40$ feet of fence for three sides of a rectangular plot; the fourth side is an
existing wall. If the two sides perpendicular to the wall are each $x$ feet long, the enclosed
area in square feet is $A(x)=x(40-2x)$.

\begin{enumerate}[leftmargin=1.6em, itemsep=6pt, topsep=4pt]
  \item Expand $A(x)$ and write it in standard form.
\begin{work}
  A(x) &= x(40-2x) \\
       &= 40x-2x^2 \\
       &= -2x^2+40x
\end{work}

  \item The factored form $A(x)=x(40-2x)$ has two zeros. Find them, then say what each one means
        for the garden.
\begin{work}
  x &= 0 \text{ or } 40-2x = 0 \\
  x &= 0 \text{ or } x = 20
\end{work}
\ansline{At $x=0$ and $x=20$ the plot collapses to a line --- no enclosed area.}

  \item A parabola is symmetric, so the largest area sits exactly halfway between the two zeros.
        Find that value of $x$ and the area it encloses.
\begin{work}
  x &= \dfrac{0+20}{2} = 10 \\
  A(10) &= 10\big(40-2(10)\big) \\
        &= 10(20) = 200
\end{work}
\end{enumerate}
\end{practicebox}

\end{document}
